 % !TEX root = ../main.tex
\section{Introduction}
\label{sec:introduction}
\glsreset{LDAP}
The \gls{LDAP} is a well-known protocol to access directory servers.
The protocol is used by various applications, including Microsoft's \gls{AD}. It serves as a way for \gls{LDAP} clients to query a backend---\eg for \gls{SSO} or to enable email clients to fetch receiver addresses or the corresponding S/MIME certificates~\cite{OndaroEtAl_MINECollectingAnalyzing_2025}.
Due to their content and widespread use of \gls{LDAP} servers in organizations, it is often intended for internal use and are not reachable from the public Internet. However, there are many legitimate cases of public-facing servers, such as the aforementioned support for (secure) email.

In recent work, we investigated the security posture of the \emph{general} \gls{LDAP} server population on the Internet \cite{landscape} with a focus on unauthorized access to sensitive data and some aspects of communication security, leaving an analysis of \emph{\gls{LDAP} hosting} to future work.

In this paper, we return to the topic with a focus on the \emph{management and hosting of \gls{LDAP} across hosting providers}. Our main contributions are: a) improved \gls{LDAP} measurements, which enable us to b) detect more \gls{LDAP} products at large scale; c) analyze the types of hosting and network segments where \gls{LDAP} is found; d) analyze the \gls{PKI} practices in the \gls{LDAP} ecosystem. We make our code and data publicly available under a permissive license\footnote{\url{github.com/internettransparency/ldap-sequel}}.

We extend our previous tooling to improve the identification of \gls{LDAP} servers by a factor of 6.5, identifying many previously unknown instances in the process. Unlike prior work, we extend our measurement and analysis to \gls{GC} instances, achieving a more comprehensive view of the ecosystem.
Our product detection can help operators assess the exposure of their networks by identifying outdated instances that an attacker could exploit.
To characterize the diverse hosting practices, we drill down to the level of the hosters' network blocks to identify and describe the network segments they use for \gls{LDAP}. We use a Pareto-inspired approach to determine the hosting behavior in segments in terms of \gls{LDAP} server software and underlying operating system.

Against this background, we then investigate management practices such as the use of strong communication security and the setup of \gls{PKI}, including attempts to deploy private \glspl{PKI}.

We find that 4.3k servers run a deprecated \gls{LDAP} product. Unlike the centralization observed in other ecosystems such as Web hosting~\cite{doan2022}, DNS~\cite{Moura2015}, and email~\cite{liuWhoGotYour2021}, we observed several hosters of \gls{LDAP} servers. We report on servers using X.509 certificates with a private key that is publicly known, posing serious risks to operators and users of these servers. Our findings enable us to provide recommendations for \gls{LDAP} operators on how to configure their systems more securely, as well as guidance for hosting operators who wish to identify poorly maintained \gls{LDAP} deployments within their network segments or develop mitigation strategies.
